\chapter{Strony i HTML}

\section{Strona jako finalny produkt}
Najwazniejszym wynikiem pracy w MoonChunk sa strony HTML. W praktyce oznacza to,
ze ostatecznym celem programu jest stworzenie jednego albo wielu dokumentow, ktore
mozna otworzyc w przegladarce, opublikowac na serwerze albo dolaczyc do wiekszego
statycznego serwisu.

W MoonChunk pojedyncza strona powstaje przez instrukcje \mc{page}.

\begin{lstlisting}[style=moonchunk,caption={Najprostsza strona}]
page "" {
  content {
    <h1>Witaj</h1>
  };
};
\end{lstlisting}

\section{Jak rozumiec trase strony}
Argument przekazywany do \mc{page} okresla trase, a w konsekwencji takze nazwe pliku
wynikowego. Warto zapamietac cztery podstawowe przypadki:

\begin{itemize}
  \item \mc{page ""} generuje \mc{index.html},
  \item \mc{page "/"} rowniez generuje \mc{index.html},
  \item \mc{page "about"} generuje \mc{about.html},
  \item \mc{page "docs/getting-started"} generuje \mc{docs/getting-started.html}.
\end{itemize}

MoonChunk pomaga tutaj uzytkownikowi. Nie musisz dopisywac rozszerzenia \mc{.html}
a ani wiodacego ukosnika. Jezyk sam zamienia trase na odpowiednia sciezke pliku.

\section{Po co to ulatwienie}
Dzieki temu zapis tras jest czystszy i bardziej zblizony do myslenia o strukturze strony,
a nie o niskopoziomowej strukturze plikow. Kiedy piszesz:

\begin{lstlisting}[style=moonchunk,caption={Trasa w stylu logicznym}]
page "blog/moj-pierwszy-wpis" {
  content {
    <h1>Moj pierwszy wpis</h1>
  };
};
\end{lstlisting}

bardziej myslisz o adresie logicznym strony niz o tym, czy trzeba dopisac \mc{.html}.
To drobiazg, ale przy duzej liczbie podstron daje bardzo wygodny efekt.

\section{Blok \mc{content}}
Wewnątrz \mc{page} najwazniejszym miejscem jest \mc{content}. To tutaj piszesz wlasciwy
HTML strony. MoonChunk traktuje ten blok jako dynamiczny fragment tresci, ktory zostanie
wstawiony do wewnetrznego szablonu dokumentu.

\begin{lstlisting}[style=moonchunk,caption={Zawartosc strony w bloku content}]
page "" {
  content {
    <section>
      <h1>Dokumentacja</h1>
      <p>Opis projektu.</p>
    </section>
  };
};
\end{lstlisting}

Wazna zasada brzmi: \important{blok \mc{content} wolno umieszczac tylko w kontekscie strony}.
Nie jest to ogolny blok do wstawiania HTML gdziekolwiek. Jego rola jest konkretna:
budowanie tresci finalnej strony.

\section{Dynamiczne wstawianie wartosci}
Jedna z najmocniejszych cech MoonChunk polega na tym, ze wewnatrz \mc{content} mozna
wstawiać wartosci zmiennych i wyniki wyrazen. Sluza do tego nawiasy klamrowe:

\begin{lstlisting}[style=moonchunk,caption={Interpolacja wartosci w HTML}]
page "" {
  let title = "MoonChunk";
  let subtitle = "Jezyk do generowania HTML";

  content {
    <h1>{title}</h1>
    <p>{subtitle}</p>
  };
};
\end{lstlisting}

MoonChunk oblicza wartosci w nawiasach i zamienia je na tekst. To bardzo wygodny sposob
na tworzenie dynamicznych naglowkow, opisow, list czy sekcji.

\section{Atrybuty HTML}
MoonChunk pozwala takze budowac dynamiczne atrybuty. Przyklad:

\begin{lstlisting}[style=moonchunk,caption={Dynamiczny atrybut HTML}]
page "" {
  let slug = "kontakt";
  let label = "Przejdz do kontaktu";

  content {
    <a href={"/" + slug}>{label}</a>
  };
};
\end{lstlisting}

W takim zapisie wartosc atrybutu \mc{href} jest obliczana dynamicznie. To przydatne
zwlaszcza wtedy, gdy generujesz wiele podobnych linkow albo sekcji na podstawie danych.

\section{Czy w \mc{content} mozna mieszac HTML i logike}
Tak, ale z umiarem. MoonChunk umozliwia kilka stylow pracy.

\subsection*{Styl prosty}
Przy malych stronach logika jest niewielka, a w \mc{content} od razu wstawia sie wartosci:

\begin{lstlisting}[style=moonchunk,caption={Prosty styl budowy strony}]
page "" {
  let title = "Start";
  content {
    <h1>{title}</h1>
  };
};
\end{lstlisting}

\subsection*{Styl uporzadkowany}
Przy wiekszej logice lepiej najpierw przygotowac dane, a dopiero potem uzyc ich w HTML-u:

\begin{lstlisting}[style=moonchunk,caption={Najpierw logika, potem HTML}]
page "" {
  let title = "Oferta";
  let intro = "Sprawdz nasze uslugi";

  content {
    <header>
      <h1>{title}</h1>
      <p>{intro}</p>
    </header>
  };
};
\end{lstlisting}

Drugi styl jest zwykle czytelniejszy w rozbudowanych projektach.

\section{Petle i warunki wewnatrz strony}
W MoonChunk treść strony moze byc skladana warunkowo i iteracyjnie. Oznacza to, ze
w ramach \mc{page} mozna uzywac \mc{if}, \mc{for} i \mc{while}, a wewnatrz nich
umieszczac \mc{content}. To bardzo mocny mechanizm.

\begin{lstlisting}[style=moonchunk,caption={Content budowany warunkowo}]
page "" {
  let isLogged: bool = true;

  if (isLogged) {
    content {
      <p>Witaj ponownie.</p>
    };
  } else {
    content {
      <p>Zaloguj sie, aby zobaczyc wiecej.</p>
    };
  };
};
\end{lstlisting}

To znaczy, ze HTML nie musi byc jedna, sztywna bryla. Moze byc skladany fragment po
fragmencie w zaleznosci od danych i warunkow.

\section{Wiecej niz jeden blok \mc{content}}
Strona moze skladac sie z wielu blokow \mc{content}. Ich zawartosc zostanie polaczona
w takiej kolejnosci, w jakiej zostaly wykonane.

\begin{lstlisting}[style=moonchunk,caption={Kilka blokow content na jednej stronie}]
page "" {
  content {
    <header><h1>Sklep</h1></header>
  };

  content {
    <section><p>Promocje dnia.</p></section>
  };

  content {
    <footer><small>Koniec strony</small></footer>
  };
};
\end{lstlisting}

To zachowanie bywa bardzo wygodne przy skladaniu strony z warunkowych fragmentow,
sekcji generowanych w petli albo czesci importowanych z innych chunkow.

\section{HTML jako wynik, nie jako cel sam w sobie}
Warto pamietac, ze MoonChunk nie probuje zastapic calego ekosystemu frontendu.
Jego zadaniem jest przygotowanie poprawnego i czytelnego HTML-a. To oznacza, ze:

\begin{itemize}
  \item style CSS mozesz dolaczyc przez metadata,
  \item skrypty JavaScript mozesz dolaczyc przez metadata,
  \item sama tresc strony jest najwazniejsza w \mc{content},
  \item MoonChunk koncentruje sie na generowaniu dokumentu, a nie na runtime w przegladarce.
\end{itemize}

\section{Praktyczna rada dotyczaca czytelnosci}
Jesli blok \mc{content} robi sie bardzo dlugi, stosuj zwykle zasady porzadku znane z HTML-a:

\begin{itemize}
  \item wciecia,
  \item logiczne grupowanie sekcji,
  \item sensowne nazwy klas,
  \item nie mieszaj zbyt wielu obliczen bezposrednio w nawiasach klamrowych.
\end{itemize}

Zamiast pisac bardzo rozbudowane wyrazenie w srodku znacznika, czesto lepiej najpierw
obliczyc wartosc w zmiennej, a potem spokojnie jej uzyc w HTML-u.
